% =============================================================================
% Chapitre 2 — Architecture Globale
% Dossier d'Architecture LockBits
% Projet Annuel ESGI 2025-2026
% =============================================================================

\chapter{Architecture Globale}

Ce chapitre présente l'architecture globale du projet LockBits. Il décrit
l'organisation en couches, les composants qui constituent la solution, leurs
interactions, les flux de données, ainsi que les choix technologiques qui ont
guide sa conception. L'objectif est de fournir une vue d'ensemble claire et
structurée du système d'information.

% =============================================================================
\section{Vue d'ensemble de l'architecture}
% =============================================================================

\subsection{Modèle architectural}

LockBits adopte une \textbf{architecture en couches} qui sépare les
responsabilités en quatre niveaux distincts :

\begin{description}
  \item[Couche présentation] \\[2pt] Interface utilisateur du portail client et
    chatbot. Elle est assurée par un site web PHP 8+ sur Apache, accessible
    depuis le navigateur via le port 8080.
  \item[Couche API et services] \\[2pt] Logique applicative exposée par le serveur
    EDR en FastAPI (Python) sur le port 8001. Cette couche traite les
    requêtes des agents de détection, orchestre les scans et gère les
    intégrations externes.
  \item[Couche données] \\[2pt] Persistance des informations. Deux technologies
    cohabitent : PostgreSQL pour les données de l'EDR, et MySQL 8.0 pour les
    données du portail client. Cette séparation permet d'isoler les
    responsabilités et d'optimiser chaque stockage pour son usage.
  \item[Couche monitoring et observabilité] \\[2pt] Collecte de métriques, de logs
    et de traces via la pile Prometheus, Grafana, Loki et Alloy. Garantit
    la visibilité sur l'etat de l'infrastructure.
\end{description}

Cette organisation en couches permet une évolution indépendante de chaque
niveau, facilite la maintenance et renforce la sécurité en limitant les
points d'accès directs aux données.

\subsection{Schéma d'architecture global}

Le schéma suivant présente l'architecture globale du système. Il montre les
quatre couches évoquées, les composants internes et leurs points d'accès
respectifs.

\begin{figure}[H]
  \centering
  \begin{tikzpicture}[
    scale=0.8,
    node distance=2.5cm and 3.5cm,
    every node/.style={font=\small},
    box/.style={
      draw=primaryblue, thick, rounded corners=3pt,
      fill=primaryblue!8, minimum width=3.8cm, minimum height=1.2cm,
      align=center
    },
    db/.style={
      draw=primaryblue!60!black, thick,
      fill=primaryblue!5, minimum width=3.2cm, minimum height=1.0cm,
      align=center, cylinder, shape border rotate=90, aspect=0.25
    },
    ext/.style={
      draw=accentorange, thick, rounded corners=0pt,
      fill=accentorange!8, minimum width=3.2cm, minimum height=1.0cm,
      align=center
    },
    mono/.style={
      draw=gray!60, thick, rounded corners=2pt,
      fill=gray!8, minimum width=2.8cm, minimum height=0.8cm,
      align=center, font=\footnotesize
    },
    layer/.style={
      draw=black!20, dashed, inner sep=6pt, rounded corners=5pt,
      fill=black!2, label={[font=\small\itshape]above:#1}
    },
    arrow/.style={->, >=stealth, thick, primaryblue!60!black},
    arrowlabel/.style={font=\footnotesize, text=primaryblue!60!black, midway, fill=white, inner sep=2pt}
  ]

    % --- Couche externe (haut) ---
    \node[ext] (internet) {Internet\\Clients Web};
    \node[ext, right=of internet] (agents) {Agents EDR};

    % --- Couche présentation / API ---
    \node[box, below=1.5cm of internet] (siteweb) {Site Web\\PHP 8+ / Apache\\\footnotesize\texttt{:8080}};
    \node[box, right=of siteweb] (edrserver) {Serveur EDR\\FastAPI / Python\\\footnotesize\texttt{:8001}};

    % --- Couche données ---
    \node[db, below=of siteweb] (mysql) {MySQL 8.0\\\footnotesize\texttt{site-db :3306}};
    \node[db, below=of edrserver] (sqlite) {SQLite\\\footnotesize(EDR data)};

    % --- Couche monitoring (bas) ---
    \node[mono, below=1.8cm of mysql] (prometheus) {Prometheus};
    \node[mono, right=of prometheus] (grafana) {Grafana};
    \node[mono, right=of grafana] (loki) {Loki + Alloy};
    \node[mono, right=of loki] (dokploy) {Dokploy};

    % --- Fleches couche présentation ---
    \draw[arrow] (internet) -- node[arrowlabel] {HTTP} (siteweb);
    \draw[arrow] (agents) -- node[arrowlabel] {HTTP / API} (edrserver);

    % --- Fleches couche API -> données ---
    \draw[arrow] (siteweb) -- node[arrowlabel] {MySQL :3306} (mysql);
    \draw[arrow] (edrserver) -- node[arrowlabel] {SQLite} (sqlite);

    % --- Interactions horizontales ---
    \draw[arrow, bend left=15] (siteweb.east) to node[arrowlabel, above] {API} (edrserver.west);
    \draw[arrow, bend left=15, accentorange!60!black] (edrserver.west) to node[arrowlabel, below] {Donnees} (siteweb.east);

    % --- Fleches monitoring ---
    \draw[arrow, gray!60] (siteweb.south) -- ++(0,-0.6) -| (prometheus.north);
    \draw[arrow, gray!60] (edrserver.south) -- ++(0,-0.6) -| (grafana.north);
    \draw[arrow, gray!60] (prometheus) -- (grafana);
    \draw[arrow, gray!60] (loki) -- (grafana);

  \end{tikzpicture}
  \caption{Schéma d'architecture globale de LockBits}
  \label{fig:architecture-globale}
\end{figure}

\subsection{Interactions entre composants}

Les composants interagissent selon les modalites suivantes :

\begin{itemize}
  \item Le \textbf{site web} communique avec la base MySQL pour gèrer les
    comptes clients, les tickets et les métriques. Les données de sécurité
    sont basees sur les tickets GLPI synchronises avec le portail.
  \item Le \textbf{serveur EDR} recoit les données des agents via une API
    REST authentifiee par jeton (\texttt{AUTH\_TOKEN}). Il les stocke en
    base PostgreSQL et peut declencher des analyses aupres de VirusTotal.
  \item Le \textbf{chatbot} intégré au site web dialogue avec l'API Claude
    d'Anthropic pour repondre aux questions des visiteurs. Il peut voir la
    discussion du ticket GLPI affichee sur la page courante pour fournir
    une assistance contextuelle.
  \item La \textbf{pile de monitoring} collecte les métriques exposées par
    chaque service via leurs endpoints \texttt{/metrics} et centralisé les
    logs via Loki et Alloy.
\end{itemize}

% =============================================================================
\section{Principes architecturaux fondamentaux}
% =============================================================================

L'architecture de LockBits repose sur six principes fondamentaux qui
guident l'ensemble des décisions de conception et de déploiement.

\subsection*{1. Modularité et découplage}

Chaque composant est développé, versionné et conteneurisé indépendamment
via des sous-modules Git. Le dépôt principal (\texttt{Lockbits-ESGI/main})
joue un rôle d'orchestration : il référence les sous-modules
\texttt{edr/} et \texttt{site\_lockbits/} sans contenir leur code source.
Cette séparation permet de :
\begin{itemize}
  \item Travailler sur chaque composant sans impact sur les autres;
  \item Réutiliser des composants dans d'autres contextes;
  \item Maintenir des cycles de livraison indépendants.
\end{itemize}

\subsection*{2. Conteneurisation systématique}

Tous les services sont empaquetés dans des images Docker multi-architecture
(\texttt{linux/amd64} et \texttt{linux/arm64}). Les images sont publiées
sur le GitHub Container Registry (GHCR) avec des tags sémantiques
(\texttt{:latest}, \texttt{:develop}, \texttt{:X.Y.Z}). Cette approche
garantit la reproductibilité des environnements du développement a la
production.

\subsection*{3. Isolation réseau}

Le déploiement utilise deux réseaux Docker distincts :
\begin{itemize}
  \item \textbf{Réseau \texttt{frontend}} : expose les ports des services
    accessibles depuis l'extérieur (EDR :8001, Site :8080).
  \item \textbf{Réseau \texttt{backend}} : isole la base de données MySQL
    du reste. Seul le site web y est connecté.
\end{itemize}
Cette séparation garantit qu'une compromission d'un service expose ne donne
pas un accès direct a la couche de données.

\subsection*{4. Observabilité}

La pile Prometheus, Grafana, Loki et Alloy fournit une visibilité complète
sur l'etat du système :
\begin{itemize}
  \item \textbf{Prometheus} collecte les métriques techniques (CPU, mémoire,
    requêtes, latence).
  \item \textbf{Grafana} offre des tableaux de bord centralisés.
  \item \textbf{Loki et Alloy} aggregent et analysent les logs de chaque
    conteneur.
\end{itemize}

\subsection*{5. Infrastructure as Code}

L'ensemble de l'infrastructure est décrit dans des fichiers versionnés :
\begin{itemize}
  \item Fichiers \texttt{docker-compose.*.yml} pour les environnements
    dev, preprod et production.
  \item Pipelines CI/CD dans \texttt{.github/workflows/} pour les tests,
    builds, scans et déploiements automatises.
  \item Scripts \texttt{install.sh}, \texttt{install-agent.sh},
    \texttt{backup.sh} et \texttt{restore.sh} pour la maintenance.
\end{itemize}

\subsection*{6. Sécurité des la conception}

La sécurité est integree a chaque etape du cycle de vie :
\begin{itemize}
  \item Scan automatique des images Docker avec Trivy dans la pipeline CI.
  \item Gestion des secrets via fichier \texttt{.env} et GitHub Secrets.
  \item Protection CORS et rate limiting sur l'API EDR.
  \item Authentification par jeton pour les agents EDR.
  \item Fichier \texttt{.gitignore} strict pour empecher les fuites de
    données sensibles.
\end{itemize}

% =============================================================================
\section{Composants majeurs et interactions}
% =============================================================================

\subsection{Stack technique}

Le tableau~\ref{tab:stack-technique} recapitule les technologies utilisees
pour chaque composant de la solution.

\begin{table}[H]
  \centering
  \caption{Stack technique de LockBits}
  \label{tab:stack-technique}
  \small
  \begin{tabular}{lccc}
    \toprule
    \textbf{Composant} & \textbf{Technologie} & \textbf{Sous-module} & \textbf{Port} \\
    \midrule
    Serveur EDR        & Python / FastAPI     & \texttt{edr/}        & 8001 \\
    Agent EDR          & Python + PyInstaller & \texttt{edr/agent/}  & --- \\
    Site web           & PHP 8+ / Apache      & \texttt{site\_lockbits/} & 8080 \\
    Base de données site & MySQL 8.0          & ---                  & 3306 \\
    Base de données EDR  & PostgreSQL         & ---                  & --- \\
    Chatbot            & PHP / Claude API     & \texttt{site\_lockbits/} & --- \\
    CI/CD              & GitHub Actions       & ---                  & --- \\
    Registry           & GHCR                 & ---                  & --- \\
    Monitoring         & Prometheus / Grafana / Loki & ---          & --- \\
    Déploiement        & Dokploy              & ---                  & --- \\
    \bottomrule
  \end{tabular}
\end{table}

\subsection{Flux de données}

Le schéma suivant illustre les flux de données entre les differents
composants du système. Chaque fleche represente un echange de données
avec son protocole et sa direction.

\begin{figure}[H]
  \centering
  \begin{tikzpicture}[
    scale=0.85,
    node distance=2.5cm and 3.5cm,
    every node/.style={font=\small},
    comp/.style={
      draw=primaryblue, thick, rounded corners=3pt,
      fill=primaryblue!6, minimum width=3.0cm, minimum height=1.0cm,
      align=center
    },
    extcomp/.style={
      draw=accentorange, thick, rounded corners=2pt,
      fill=accentorange!6, minimum width=2.8cm, minimum height=0.9cm,
      align=center
    },
    db/.style={
      draw=primaryblue!60!black, thick,
      fill=primaryblue!4, minimum width=2.8cm, minimum height=0.9cm,
      align=center, cylinder, shape border rotate=90, aspect=0.25
    },
    arrow/.style={->, >=stealth, thick, primaryblue!50!black},
    arrowrev/.style={<-, >=stealth, thick, primaryblue!50!black},
    arrowboth/.style={<->, >=stealth, thick, primaryblue!50!black},
    label/.style={font=\footnotesize, text=primaryblue!60!black, fill=white, inner sep=2pt, align=center}
  ]

    % --- Ligne du haut : composants externes ---
    \node[extcomp] (agent)      {Agent EDR};
    \node[extcomp, right=of agent] (vt)       {VirusTotal};
    \node[extcomp, right=of vt] (glpi)     {GLPI};
    \node[extcomp, right=of glpi] (claude)   {Claude API};

    % --- Ligne du milieu : serveurs ---
    \node[comp, below=of agent] (edr)     {Serveur EDR\\FastAPI};
    \node[comp, right=of edr] (site)  {Site Web\\PHP / Apache};

    % --- Ligne du bas : bases de données ---
    \node[db, below=of edr] (sqldb)   {SQLite\\EDR data};
    \node[db, right=of sqldb] (mysqldb) {MySQL 8.0\\site-db};

    % --- Chatbot intégré au site ---
    \node[comp, right=4.5cm of site, yshift=-1.1cm] (chatbot) {Chatbot\\PHP};

    % --- Fleches Agent EDR -> EDR Server ---
    \draw[arrow] (agent) -- node[label] {Heartbeat\\+ Evenements} (edr);
    \draw[arrowrev] (edr) -- node[label] {Reponses\\+ Commandes} (agent);

    % --- EDR Server -> VirusTotal ---
    \draw[arrow] (edr.north) -- ++(0,0.6) -| node[label, above] {Scan fichiers\\suspects} (vt.south);

    % --- EDR Server -> GLPI ---
    \draw[arrow] (edr.north) -- ++(0,0.6) -| node[label, above] {Creation tickets\\OAuth2} (glpi.south);

    % --- EDR Server -> SQLite ---
    \draw[arrow] (edr) -- node[label] {Stockage} (sqldb);
    \draw[arrowrev] (sqldb) -- node[label] {Lecture} (edr);

    % --- Site Web -> MySQL ---
    \draw[arrowboth] (site) -- node[label] {Donnees client\\+ Tickets} (mysqldb);

    % --- Site Web -> GLPI SSO ---
    \draw[arrowboth] (site.north) -- ++(0,0.8) -| node[label, above] {SSO OAuth\\Synchronisation} (glpi.south);

    % --- Site Web <-> EDR Server ---
    \draw[arrowboth, bend left=20] (site.west) to node[label, above] {API REST\\Donnees sécurité} (edr.east);

    % --- Chatbot -> Claude API ---
    \draw[arrow] (chatbot) -- node[label] {Questions\\utilisateur} (claude);
    \draw[arrowrev] (claude) -- node[label] {Reponses} (chatbot);

  \end{tikzpicture}
  \caption{Flux de données entre les composants}
  \label{fig:flux-données}
\end{figure}

\subsection{Organisation en sous-modules Git}

Le projet suit une architecture multi-dépôts orchestree par un dépôt
principal. Chaque composant applicatif possede son propre dépôt Git, ce qui
permet un développement parallele et un versionnement indépendant.

\subsubsection{Structure des dépôts}

\begin{verbatim}
Lockbits-ESGI/
├── main/                    # Dépôt d'orchestration (docker-compose, CI/CD, scripts)
│   ├── edr/                 # Sous-module -> Lockbits-ESGI/edr
│   ├── site_lockbits/       # Sous-module -> Lockbits-ESGI/site_lockbits
│   ├── docker-compose.yml
│   ├── docker-compose.prod.yml
│   ├── .github/workflows/
│   └── install.sh
│
├── edr/                     # Dépôt autonome
│   ├── server/              # Serveur FastAPI
│   ├── agent/               # Agent EDR Python
│   └── Dockerfile
│
└── site_lockbits/           # Dépôt autonome
    ├── client/              # Application PHP
    ├── chatbot/             # Chatbot Claude API
    └── Dockerfile
\end{verbatim}

\subsubsection{Mise a jour automatique des sous-modules}

Un workflow GitHub Actions (\texttt{auto-update-submodules.yml}) s'execute
toutes les heures. Il verifie si les sous-modules \texttt{edr} et
\texttt{site\_lockbits} pointent vers le dernier commit de leur branche par
defaut. Si une mise a jour est disponible, le workflow :

\begin{enumerate}
  \item Met a jour la référence du sous-module avec
    \texttt{git submodule update --remote};
  \item Cree un commit de mise a jour;
  \item Pousse sur \texttt{main}, ce qui declenche la pipeline CI/CD de
    reconstruction des images.
\end{enumerate}

Cette automatisation garantit que les images Docker publiées sur GHCR
refletent toujours la derniere version du code source, sans intervention
manuelle.

\subsubsection{Mise a jour manuelle}

Si un besoin immediat se présente, la commande suivante permet de forcer la
mise a jour depuis le serveur :

\begin{lstlisting}[language=bash, caption={Mise a jour manuelle des sous-modules}]
bash install.sh --update-sources
\end{lstlisting}

Cette commande execute les memes operations que le workflow automatique :
mise a jour des références, commit et push.

% =============================================================================
\section{Justification des choix technologiques}
% =============================================================================

Chaque technologie a été choisie en fonction des contraintes du projet :
equipe reduite, delai de huit mois, objectifs pedagogiques et realisme
professionnel. Cette section detaillle les motivations derriere chaque
choix.

\subsection{FastAPI pour le serveur EDR}

Le choix de FastAPI (Python) pour le serveur backend repose sur plusieurs
avantages decisifs :

\begin{itemize}
  \item \textbf{Performances asynchrones} : base sur ASGI (uvicorn),
    FastAPI gère nativement les requêtes concurrentes, ce qui est
    essentiel pour un serveur traitant des données en temps reel depuis de
    multiples agents.
  \item \textbf{Documentation automatique} : l'interface \texttt{/docs}
    genere une documentation interactive OpenAPI sans configuration
    additionnelle, facilitant le debogage et l'intégration.
  \item \textbf{Validation des données} : via Pydantic, les schémas de
    données sont declares en Python et valides automatiquement, reduisant
    les risques d'erreurs.
  \item \textbf{Legèreté} : démarrage rapide, empreinte mémoire reduite,
    ideal pour un conteneur Docker.
\end{itemize}

\textbf{Alternatives ecartees} : Flask manque de support asynchrone natif
et impose des extensions tierces pour la validation. Django est trop lourd
pour un service d'API specialise comme l'EDR.

\subsection{PHP / Apache pour le portail client}

Le portail client utilise PHP 8+ sur Apache, un choix guide par la
simplicite et l'adequation au besoin :

\begin{itemize}
  \item \textbf{Déploiement simple} : la stack LAMP est mature et ne
    necessite pas de build frontend. Le code PHP est directement executable
    sur le serveur.
  \item \textbf{Pages serveur} : le portail affiche des données issues de
    MySQL (tableau de bord, tickets, historique). PHP excelle dans ce modèle
    de rendu serveur.
  \item \textbf{Compatibilite MySQL} : les pilotes natifs PHP pour MySQL
    sont performants et bien documentes.
  \item \textbf{Couts de maintenance} : une equipe reduite peut maintenir
    un site PHP sans infrastructure frontend complexe.
\end{itemize}

\textbf{Alternatives ecartees} : React/Next.js auraient necessite un
serveur Node.js supplémentaire et un build frontend, ajoutant de la
complexite sans benefice significatif pour un portail client a pages
serveur.

\subsection{MySQL 8.0 pour la base de données du site}

\begin{itemize}
  \item \textbf{Fiabilité} : MySQL est eprouve depuis des decennies dans
    des environnements de production varies.
  \item \textbf{Compatibilité PHP} : l'ecosysteme PHP est historiquement
    lie a MySQL, avec des pilotes matures et optimises.
  \item \textbf{Sauvegardes simples} : \texttt{mysqldump} permet des
    sauvegardes coherentes sans outils tiers.
  \item \textbf{Performance} : pour le volume de données attendu (clients,
    tickets, métriques), MySQL 8.0 offre des performances amplement
    suffisantes.
\end{itemize}

\subsection{SQLite pour la base de données EDR}

\begin{itemize}
  \item \textbf{Simplicité} : pas de serveur de base de données a gèrer.
    Le fichier SQLite est stocke dans un volume Docker.
  \item \textbf{Volume adapté} : les données EDR (événements, alertes,
    scans) restent moderees pour une PME.
  \item \textbf{Migration possible} : l'architecture permet de basculer
    vers PostgreSQL en production sans modification majeure du code,
    via la configuration \texttt{DATABASE\_URL}.
\end{itemize}

\subsection{Docker Compose pour l'orchestration}

\begin{itemize}
  \item \textbf{Simplicité} : Docker Compose definit l'ensemble des
    services, réseaux et volumes dans un fichier YAML lisible.
  \item \textbf{Un seul hôte} : le déploiement cible un serveur unique
    (VM Proxmox), ce qui rend Kubernetes surdimensionne.
  \item \textbf{Realisme} : Docker Compose est utilise en production par
    de nombreuses PME et startups pour des charges modestes.
  \item \textbf{Compatible Dokploy} : la plateforme de déploiement
    utilisee supporte nativement les stacks Docker Compose.
\end{itemize}

\textbf{Alternative e cartée} : Kubernetes (K3s) a été ecarte car trop
complexe pour l'infrastructure actuelle, mais reste une évolution viable
et coherente a plus long terme.

\subsection{GitHub Actions pour le CI/CD}

\begin{itemize}
  \item \textbf{Intégration native} : le projet est heberge sur GitHub.
    GitHub Actions s'intégré sans configuration tierce.
  \item \textbf{Zéro infrastructure} : les runners sont fournis par GitHub,
    sans serveur CI a gèrer.
  \item \textbf{Workflows homemade} : les pipelines sont entierement
    écrits sur mesure (faits maison), sans utiliser d'actions
    pre-construites pour la logique metier.
  \item \textbf{Gratuit} : pour un projet open-source, les quotas gratuits
    sont largement suffisants.
\end{itemize}

\subsection{GHCR pour le registry Docker}

\begin{itemize}
  \item \textbf{Gratuit pour les organisations GitHub} : les images publiques
    sont autorisees sans limite de stockage.
  \item \textbf{Intégration GitHub} : pas de compte Docker Hub sépare a
    gèrer. L'authentification utilise le token GitHub existant.
  \item \textbf{Limites de pull élevées} : GHCR offre des quotas plus
    genereux que Docker Hub pour les images publiques.
  \item \textbf{Scan de vulnérabilités} : GHCR propose des fonctionnalites
    de scan de sécurité, qui peuvent être utilisees pour analyser les
    images stockees.
\end{itemize}

% =============================================================================
\section{Évolution de l'architecture}
% =============================================================================

L'architecture de LockBits a été concue pour evoluer par phases. Cette
section décrit la roadmap d'implementation, du MVP jusqu'aux extensions
futures.

\subsection{Phase MVP -- Fondations}

La phase initiale etablit la base fonctionnelle du système :

\begin{itemize}
  \item Déploiement Docker Compose sur un hote unique avec l'ensemble de
    la stack.
  \item Serveur EDR fonctionnel : API REST, base SQLite, heartbeat des
    agents, scan local.
  \item Site web operationnel : authentification, tableau de bord, tickets.
  \item Chatbot intégré au site avec l'API Claude.
  \item Pipeline CI de base : lint, build, push vers GHCR.
\end{itemize}

\subsection{Phase Consolidation -- Mise en production}

Cette phase renforce la robustesse et la maintenabilite :

\begin{itemize}
  \item CI/CD complète avec tests unitaires, intégration et scan Trivy.
  \item Déploiement automatise via Dokploy.
  \item Monitoring complète : Prometheus, Grafana, Loki, Alloy.
  \item Scripts de sauvegarde et restauration automatises.
  \item Hardening de la sécurité : CORS, rate limiting, secrets.
  \item Documentation operationnelle et d'architecture.
\end{itemize}

\subsection{Phase Extension -- Passage à l'échelle}

Les évolutions futures, détaillées dans le chapitre~\ref{ch:extensions},
incluent :

\begin{itemize}
  \item Migration vers Kubernetes (K3s) pour la haute disponibilité.
  \item SIEM complète avec Wazuh ou ELK.
  \item Authentification OAuth2 complète (Google, Microsoft).
  \item Interface d'administration web dédiée a l'EDR.
  \item Tests de pénétration pourraient être envisages via des outils
    comme OWASP ZAP intégrés a la pipeline CI/CD.
  \item Support IPv6.
\end{itemize}

Cette approche par phases permet de livrer rapidement un produit
fonctionnel tout en conservant une trajectoire claire vers une
infrastructure professionnelle et résiliente.

% =============================================================================
\endinput
% =============================================================================
